% Imporditud: tools/import_eksperimendid.py
% Allikas: Eesti füüsikaolümpiaadi eksperimendi ülesannete
%   kogu (Rael Kalda, 2023), ülesanne Ü157, lahendus L157.
\setAuthor{Jaan Kalda}
\setRound{piirkonnavoor}
\setYear{2021}
\setNumber{G E2}
\setDifficulty{4}
\setTopic{dünaamika}
\setVahendid{joonlaud, liitrine anum veega, veeklaas vee valamiseks, mille seinale on tehtud märge umbes 150 mL veehulgale vastavale nivookõrgusele, tühi liitrine õhukesest kartongist piima või mahalapakk, mille ülemine tahk on ära lõigatud (massi vähendamiseks ja vee sissevalamise hõlbustamiseks), mõned peenikeste spagettide kõrred (Nr 5 või Nr 7), majapidamispaber kuivatamiseks}
\setPunktid{12}
\setAllikas{Eksperimendi kogu Ü157, lahendus lk 118}
\setLahendusseis{toores}

\prob{Spagetikõrs}
Spagetikõrrega lükates ei ole võimalik avaldada piiramatult suurt jõudu: isegi, kui kõrs oleks absoluutselt sirge ja lükata täpselt piki kõrt, siis teatud jõust alates muutub kõrs ebastabiilseks ja paindub. Eeldades, et maksimaalne avaldatav jõud sõltub fikseeritud kõrre jämeduse juures kõrre pikkusest L vastavalt valemile Fm = AL$\alpha$ (kus A ja $\alpha$ on konstandid), määrake eksperimentaalselt astmenäitaja $\alpha$ väärtus.

\hint

\solu
Valame piimapakki vett - alustame ühest kopsikutäiest ning lõpetame umbes ühe liitriga. Iga vee koguse juures üritame lükata piimapakki spagetikõrrega mööda lauapinda - leiame sellise maksimaalse sõrmega kinnihoidmiskoha kauguse kõrre otsast, mille puhul on veel võimalik lükata. Hoolitseme, et paki põhi ja laud oleksid kogu aeg kuivad, sest vesi võib muuta hõõrdetegurit. Kanname graafiku vertikaalteljele anumas oleva vee massi (ühikuteks kasutame kopsikutäit) logaritmi ja horisontaalteljele - kõrre pikkuse logaritmi. Sirge tõusunurga tangens annab astmenäitaja. Alternatiivina võib teha oletusi astmenäitaja väärtuse kohta ja katse-eksituse meetodil leida selline astmenäitaja, mille korral teljestikus m - L$\alpha$ moodustub andmepunktidest sirgjoon. Viimase meetodi eeliseks on see, et tühja anuma mass ei põhjusta kõrvalekaldeid sirgjoonelisest graafikust. Tulemuseks leiame $\alpha \approx -$2.

Hindamine Idee kasutada osaliselt veega täidetud piimapaki ja laua vahelist hõõrdejõudu reguleeritava väärtusega jõu saamiseks (1 p). Idee lükata kõrrega leides kõrrest kinni hoidmiseks selline maksimaalne kaugus piimapakist, mille korral on veel võimalik pakki nihutada (2 p). Kogutud niimitu andmepunkti, kui antud mõõteanum võimaldab: liitrise piimapaki ja 150 mL mõõtenõu juures ühest kuni kuue kopsikutäieni, kokku kuus andmepunkti: 3 p. Kui on viis andmepunkti, kuid kirjutatud, et kuues kopsikutäis ei mahtunud enam ära, siis anda ka viie andmepunkti eest 3 p. Vastasel korral 5 andmepunkti eest 2,5 p. Osalised punktid: 4 andmepunkti: 2 p; 3 andmepunkti: 1 p; 2 andmepunkti: 0,5 p. Iga arvesse mineva andmepunkti jaoks peab olema näidatud vee kogus kopsikute arvuna ning spagetikõrre lükkava osa pikkus. Kui pikkusühikul pole näidatud ühikut, siis lahutatakse kogusummast 0,5 p. Valitud sobilik graafiline meetod: kas log-log teljestik või proovitud katse-eksituse meetodil m-L$\alpha$ teljestikku erinevate $\alpha$ väärtustega (3 p). Osalised punktid: kui astmenäitaja pole leitud graafiliselt vaid on arvutatud andmepunktipaaride põhjal, siis iga erineva andmepunktipaari põhjal arvutatud väärtus annab 0,5 p, kuni maksimaalselt 2 p (mis antakse siis, kui astmenäitaja on leitud nelja erineva andmepunkti paari jaoks). Mõõtmise täpsuse eest: kui leitud astmenäitaja on vahemikus -2,3 kuni -1,7, siis 2 p, kui tulemus on sellest vahemikust väljas, kuid jääb vahemikku -1,5 kuni -2,5, siis 1 p. Täiendavate andmepunktide leidmine sel teel, et lükatakse suurema veekoguse korral kahe või enama kõrrega korraga: 1 p.
\probend
